% Generated draft: requires mathematical and source review.
\documentclass[11pt]{article}
\usepackage[margin=25mm]{geometry}
\usepackage{fontspec}
\setmainfont{Latin Modern Roman}
\usepackage{amsmath,amssymb,mathtools,unicode-math}
\setmathfont{Latin Modern Math}
\usepackage{enumitem,needspace,xurl,hyperref}
\setlength{\parindent}{0pt}
\setlength{\parskip}{5pt}
\setlength{\emergencystretch}{3em}
\newcommand{\N}{\mathbb{N}}
\newcommand{\Z}{\mathbb{Z}}
\newcommand{\Q}{\mathbb{Q}}
\newcommand{\R}{\mathbb{R}}
\newcommand{\C}{\mathbb{C}}
\newcommand{\F}{\mathbb{F}}
\newcommand{\A}{\mathbb{A}}
\newcommand{\PP}{\mathbb{P}}
\newcommand{\E}{\mathbb{E}}
\newcommand{\eps}{\varepsilon}
\newcommand{\dd}{\,\mathrm{d}}
\begin{document}
\section*{UnsolvedMath ID: 30000089}
\textbf{Source problem code:} catalogue entry
\subsection{Definitions and mathematical statement}
The corrected-and-verified target is: for a saturated weighted digraph $G$, does its associated ideal $J_G$ have a Scarf-type minimal free resolution? The quantifier is over the intended family of such digraphs and the ideal construction in the report.

A free resolution of a quotient $S/J_G$ is an exact complex
\[
 \cdots\longrightarrow F_2\longrightarrow F_1\longrightarrow F_0
 \longrightarrow S/J_G\longrightarrow0
\]
of free modules over the relevant polynomial ring $S$. In a graded setting it is minimal when differential matrices have entries in the homogeneous maximal ideal; equivalently, no cancellable free summands remain. Resolving $J_G$ itself uses the corresponding shifted convention, which must be fixed consistently.

For comparison, if an ideal is monomial with minimal generators $u_1,\ldots,u_r$, its ordinary Scarf complex consists of generator subsets whose least-common-multiple labels occur for no other subset. Saying that it supports a minimal free resolution requires exactness of the associated labeled complex, not merely constructing that simplicial complex. The term ``Scarf-type'' in the question may refer to an analogous construction for the report's graph ideal. The excerpt does not supply the definition of saturated weighted digraph, its admissible weights, the generators of $J_G$, the base field or grading, or the precise meaning of Scarf-type. It therefore does not justify assuming that $J_G$ is a monomial, edge, or lattice ideal or substituting one of those constructions.

The attached literature assessment checked on 21 August 2026 describes a different-looking fragment saying only that an unspecified conjecture is known for $a$-stable ideals, with its conjecture and quantifiers missing. That assessment does not match the corrected weighted-digraph question closely enough to establish its status or supply its definitions. The $a$-stable citation is retained as supplied provenance, not treated as a theorem about $J_G$. The registry's open label remains unverified pending reconciliation with the original report.
\subsection{Short English statement}
Does the ideal attached to every intended saturated weighted digraph admit a Scarf-type minimal resolution? The graph-ideal definitions are missing, and the supplied literature assessment discusses an unmatched $a$-stable-ideal fragment, so it cannot certify this question's status.
\subsection{Sources}
\begin{itemize}
\item \emph{Combinatorial Commutative Algebra} (2005), Oberwolfach Reports, DOI 10.4171/owr/2004/32, pp. 1703--1746, \url{https://doi.org/10.4171/owr/2004/32}; proposed year supplied as 2004.
\item Supplied literature assessment checked 21 August 2026: Gasharov--Hibi--Peeva, ``a-stable ideals,'' \url{https://doi.org/10.1016/S0021-8693(02)00083-2}; ``Original Oberwolfach report,'' \url{https://doi.org/10.4171/owr/2004/32}. The assessment's stated fragment is not identified with the corrected graph-ideal target.
\item Local UnsolvedMath registry, permanent ID $30000089$, exact source code \texttt{catalogue entry}, created 25 August 2026 and updated 16 September 2026. Its original and clean statements both ask the weighted-digraph Scarf-resolution question; supplied verification labels it corrected and checked against the report.
\end{itemize}
\end{document}
